\documentclass[a4paper,12pt,fontset=ubuntu]{ctexart}
\usepackage[margin=2.5cm]{geometry}
\usepackage{graphicx}
\usepackage{booktabs}
\usepackage{longtable}
\usepackage{enumitem}
\usepackage{amsmath}
\usepackage{amssymb}
\usepackage{listings}
\usepackage{xcolor}
\usepackage{hyperref}
\hypersetup{colorlinks=true,linkcolor=black}

\lstset{
  basicstyle=\small\ttfamily,
  keywordstyle=\color{blue},
  commentstyle=\color{gray},
  breaklines=true,
  frame=single,
  language=Verilog
}

\title{\textbf{自适应多粒度稀疏编码架构\\AI芯片IP核专利报告书}}
\author{}
\date{\today}

\begin{document}
\maketitle
\thispagestyle{empty}
\newpage

% ============================================================================
% 第一部分：现有技术及技术问题
% ============================================================================
\section{现有技术及技术问题}

\subsection{现有技术概况}

在AI芯片加速领域，稀疏计算是降低计算量和访存开销的关键技术手段。
目前业界主流的稀疏加速方案可分为以下几类：

\subsubsection{固定粒度稀疏加速器}

以NVIDIA Ampere架构为代表，采用2:4结构化稀疏模式，
即每连续4个元素中恰好有2个非零值。该方案通过固定粒度的稀疏编码，
在硬件层面实现2倍的计算加速。

AMD CDNA架构同样采用固定稀疏模式，支持结构化稀疏矩阵运算，
但粒度选择同样受限于预设模式。

学术界的EIE（Efficient Inference Engine）采用非结构化元素级稀疏，
通过CSC（Compressed Sparse Column）格式编码权重，
支持任意稀疏模式但仅限于单一粒度。

SCNN（Sparse CNN Accelerator）采用坐标列表格式存储稀疏数据，
同样仅支持元素级单一粒度。

\subsubsection{编译时静态稀疏方案}

Google TPU和华为昇腾系列中的稀疏加速模块，
均在编译时或训练时确定稀疏模式，运行时无法根据数据分布动态调整。

MIT的SparseY虽然支持多种稀疏格式，
但格式选择在编译时完成，不具备运行时自适应能力。

\subsubsection{位图编码方案}

Sanger等方案采用位图（bitmap）标记非零元素位置，
但位图粒度固定，当稀疏率较低时位图本身的存储开销反而超过计算节省。

\subsection{现有技术存在的缺陷}

经分析，现有技术主要存在以下五个方面的技术问题：

\subsubsection{问题一：稀疏粒度单一，无法适配混合稀疏模式}

\begin{table}[htbp]
\centering
\caption{不同AI模型层的稀疏模式差异}
\begin{tabular}{llcc}
\toprule
\textbf{模型层} & \textbf{典型稀疏模式} & \textbf{最优粒度} & \textbf{现有方案粒度} \\
\midrule
Attention层（头剪枝） & 通道级 & L3 & 仅L2 \\
MLP层（2:4剪枝） & 块级 & L2 & 仅L2 \\
深层卷积（非结构化） & 元素级 & L1 & 仅L2 \\
Embedding层（低稀疏） & 稠密回退 & dense & 仅L2 \\
\bottomrule
\end{tabular}
\end{table}

如表1所示，不同网络层具有截然不同的稀疏分布特征。
现有方案仅支持单一粒度（如NVIDIA仅支持2:4块级），
导致非最优粒度的层要么无法利用稀疏加速，要么加速效率大幅下降。

\subsubsection{问题二：静态编码决策，无法适应运行时数据分布变化}

现有方案的稀疏编码策略在编译时或训练时确定，运行时无法动态调整。
然而，实际推理场景中数据分布存在显著变化：

\begin{itemize}
  \item 不同batch的输入数据稀疏度可能波动30\%以上；
  \item 动态剪枝（runtime pruning）产生的稀疏模式在推理过程中不断变化；
  \item 多任务场景下不同任务激活不同的稀疏区域。
\end{itemize}

静态策略无法应对上述变化，导致稀疏收益随数据波动大幅下降。

\subsubsection{问题三：稀疏索引开销抵消计算节省}

以位图编码为例，对于4$\times$4块级稀疏，位图占16bit；
当块内非零元素较多时，位图开销占非零数据量的比例过高：

\begin{equation}
\eta_{\text{overhead}} = \frac{N_{\text{bitmap}}}{N_{\text{nz}} \times W_{\text{data}}}
\end{equation}

其中$N_{\text{bitmap}}$为位图位数，$N_{\text{nz}}$为非零元素数，
$W_{\text{data}}$为数据位宽。
当$N_{\text{nz}} \geq 12$时，$\eta_{\text{overhead}} > 25\%$，
编码收益已不明显。

\subsubsection{问题四：低稀疏度场景缺乏回退机制}

当稀疏率低于50\%时，编码索引开销加上稀疏计算的控制开销
可能超过直接稠密计算的开销。现有方案（如2:4稀疏）
强制要求50\%稀疏率，无法在低稀疏度场景自动回退到稠密模式，
导致低稀疏度时性能反而下降。

\subsubsection{问题五：IP核复用困难}

不同算法（CNN、Transformer、LLM）的稀疏模式差异大，
现有方案难以用同一IP核覆盖。通常需要为不同算法设计不同的稀疏加速单元，
增加了芯片面积和设计复杂度。

% ============================================================================
% 第二部分：发明方案具体介绍
% ============================================================================
\section{发明方案具体介绍}

\subsection{发明概述}

本发明提出一种自适应多粒度稀疏编码架构，作为AI芯片的IP核，
能够根据运行时输入数据的稀疏特性，自动选择最优编码粒度，
实现编码效率最大化。核心创新包括：

\begin{enumerate}
  \item 定义四级稀疏粒度模型（L0位级/L1元素级/L2块级/L3通道级），支持混合粒度编码；
  \item 运行时自适应粒度决策引擎，根据稀疏率和空间聚集度动态选择粒度；
  \item 统一混合编码格式，支持不同粒度在统一框架下编码和解码；
  \item 零开销稠密回退机制，在低稀疏度场景自动切换为稠密计算。
\end{enumerate}

\subsection{多粒度稀疏模型}

\subsubsection{粒度级别定义}

本发明定义四级稀疏粒度，从细到粗：

\begin{table}[htbp]
\centering
\caption{四级稀疏粒度定义}
\begin{tabular}{cllll}
\toprule
\textbf{级别} & \textbf{名称} & \textbf{编码单位} & \textbf{索引格式} & \textbf{适用场景} \\
\midrule
L0 & 位级稀疏 & 单个bit & 位掩码 & 量化后零值bit \\
L1 & 元素级稀疏 & 单个权重/激活 & (row,col)坐标 & 非结构化剪枝 \\
L2 & 块级稀疏 & $N \times M$块 & 块位图 & 半结构化稀疏(2:4) \\
L3 & 通道级稀疏 & 整行/列/通道 & 通道掩码 & 结构化剪枝 \\
\bottomrule
\end{tabular}
\end{table}

\textbf{关键特征}：同一数据流的不同区域可使用不同粒度编码，
由硬件在运行时动态切换，粒度选择由2bit标识位编码在统一格式中。

\subsubsection{各粒度编码格式}

\paragraph{L1元素级编码}

采用坐标列表（COO）格式，每个非零元素记录其坐标和值：

\begin{equation}
\text{L1}: \{(r_i, c_i, v_i)\}_{i=0}^{N_{\text{nz}}-1}
\end{equation}

存储开销：$S_{L1} = N_{\text{nz}} \times (2W_{\text{idx}} + W_{\text{data}})$

适用于高稀疏率（$>$80\%）且空间分散的场景。

\paragraph{L2块级编码}

将张量划分为$B_H \times B_W$的块，每块用位图标记非零位置：

\begin{equation}
\text{L2}: \{\text{bitmap}_j, \{v_k\}_{k \in \text{nz}(j)}\}_{j=0}^{N_{\text{blk}}-1}
\end{equation}

存储开销：$S_{L2} = N_{\text{blk}} \times B_H B_W + N_{\text{nz}} \times W_{\text{data}}$

适用于中等稀疏率（50\%--80\%）且具有一定空间聚集性的场景。

\paragraph{L3通道级编码}

以通道为最小编码单位，用掩码标记非零通道：

\begin{equation}
\text{L3}: \{\text{ch\_mask}, \{V_c\}_{c \in \text{nz\_ch}}\}
\end{equation}

存储开销：$S_{L3} = N_{\text{ch}} + N_{\text{nz\_ch}} \times H \times W \times W_{\text{data}}$

适用于低稀疏率（30\%--50\%）或通道级结构化剪枝场景。

\paragraph{稠密回退}

当稀疏率$<$30\%时，编码索引开销超过计算节省，
直接存储原始稠密数据，零额外开销：

\begin{equation}
S_{\text{dense}} = N_{\text{total}} \times W_{\text{data}}
\end{equation}

\subsection{统一混合编码格式}

本发明设计的统一编码格式头部标识粒度，载荷按粒度编码：

\begin{table}[htbp]
\centering
\caption{统一编码包格式}
\begin{tabular}{lll}
\toprule
\textbf{字段} & \textbf{位宽} & \textbf{说明} \\
\midrule
granularity   & 2bit   & 粒度标识：00=L1, 01=L2, 10=L3, 11=dense \\
shape\_info   & 46bit  & 张量形状信息(rows + cols + channels) \\
nz\_count     & 8bit   & 非零元素计数 \\
index\_len    & 8bit   & 索引数据有效长度 \\
index\_data   & 128bit & 索引/位图数据（粒度相关） \\
value\_data   & 变长   & 非零值紧凑排列 \\
\bottomrule
\end{tabular}
\end{table}

\textbf{核心优势}：不同通道、不同层可使用不同粒度标识，
解码器根据2bit粒度标识分发至对应解码通路，实现零开销粒度切换。

\subsection{自适应粒度决策引擎}

\subsubsection{整体架构}

自适应粒度决策引擎由三部分组成：

\begin{enumerate}
  \item \textbf{稀疏感知分析器}：运行时统计输入张量的稀疏率和空间聚集度；
  \item \textbf{粒度决策状态机}：根据统计结果从多级粒度中自适应选择；
  \item \textbf{可编程配置表}：支持通过寄存器配置阈值参数。
\end{enumerate}

\subsubsection{稀疏感知分析器}

稀疏感知分析器在每个张量处理周期内，实时统计以下两个关键指标：

\paragraph{稀疏率}定义为零元素占总元素的比例：

\begin{equation}
\rho_{\text{sparse}} = 1 - \frac{N_{\text{nz}}}{N_{\text{total}}}
\end{equation}

硬件实现中采用Q8.8定点数格式表示，通过逐次逼近除法器计算，
避免硬件除法器的高面积开销。

\paragraph{空间聚集度}定义为非零块数与理论最大非零块数的比值：

\begin{equation}
\rho_{\text{cluster}} = \frac{N_{\text{nz\_blk}}}{N_{\text{nz}} / (B_H \times B_W)}
\end{equation}

聚集度越高，说明非零元素在空间上越集中，更适合块级编码；
聚集度越低，说明非零元素越分散，更适合元素级编码。

\subsubsection{决策算法}

决策引擎采用阈值比较加启发式规则的方式，实现低延迟决策：

\begin{table}[htbp]
\centering
\caption{粒度决策规则}
\begin{tabular}{llll}
\toprule
\textbf{条件} & \textbf{选择粒度} & \textbf{原因编码} & \textbf{编码格式} \\
\midrule
$\rho_{\text{sparse}} \geq T_{L1}$ 且 $\rho_{\text{cluster}} < T_c$ & L1元素级 & 0x01 & 坐标列表 \\
$\rho_{\text{sparse}} \geq T_{L1}$ 且 $\rho_{\text{cluster}} \geq T_c$ & L2块级 & 0x02 & 块位图 \\
$T_{L2} \leq \rho_{\text{sparse}} < T_{L1}$ & L2块级 & 0x03 & 块位图 \\
$T_{L3} \leq \rho_{\text{sparse}} < T_{L2}$ & L3通道级 & 0x04 & 通道掩码 \\
$\rho_{\text{sparse}} < T_{L3}$ & 稠密回退 & 0x05 & 原始数据 \\
\bottomrule
\end{tabular}
\end{table}

默认阈值配置：$T_{L1} = 0.80$，$T_{L2} = 0.50$，$T_{L3} = 0.30$，$T_c = 0.50$。
所有阈值通过AXI-Lite寄存器可编程配置。

\subsubsection{运行时自适应机制}

决策引擎支持运行时动态更新粒度选择：

\begin{enumerate}
  \item 每个张量处理周期（由tensor\_last信号标识）自动重新评估；
  \item 可配置更新周期：每$N$个batch重新评估一次（$N$通过寄存器配置）；
  \item 支持强制粒度模式：通过cfg\_force\_en寄存器强制指定粒度，用于调试。
\end{enumerate}

\subsection{硬件实现架构}

\subsubsection{顶层模块划分}

IP核的RTL实现包含以下关键模块：

\begin{table}[htbp]
\centering
\caption{IP核模块列表}
\begin{tabular}{llp{7cm}}
\toprule
\textbf{模块名} & \textbf{功能} & \textbf{关键实现} \\
\midrule
sparsity\_analyzer   & 稀疏感知分析 & $countones$零值检测 + 逐次逼近除法器 \\
granularity\_decision & 粒度决策引擎 & 阈值比较 + 状态机 \\
encoder\_l1          & L1元素级编码 & COO坐标生成 \\
encoder\_l2          & L2块级编码 & 位图生成 + 非零值紧凑排列 \\
encoder\_l3          & L3通道级编码 & 通道掩码 + 通道数据紧凑排列 \\
decoder\_l1          & L1元素级译码 & 坐标写入 \\
decoder\_l2          & L2块级译码 & 位图扩展恢复 \\
decoder\_l3          & L3通道级译码 & 掩码扩展恢复 \\
granularity\_router  & 粒度路由器 & 1-to-4 demux \\
encoded\_packet\_formatter & 编码包格式化 & 统一格式打包 \\
\bottomrule
\end{tabular}
\end{table}

\subsubsection{数据通路}

完整编码数据通路为：

\begin{enumerate}
  \item 输入张量通过AXI-Stream接口进入稀疏感知分析器；
  \item 分析器输出稀疏率和聚集度统计至决策引擎；
  \item 决策引擎选择最优粒度，控制粒度路由器分发数据；
  \item 对应编码器通路完成编码，打包为统一格式输出；
  \item 译码时根据包中粒度标识，路由至对应译码通路恢复数据。
\end{enumerate}

\subsubsection{接口定义}

IP核对外提供两类接口：

\paragraph{AXI-Lite配置接口}用于寄存器读写，配置阈值、块大小、更新周期等参数。

\paragraph{AXI-Stream数据接口}用于张量数据流输入和编码数据输出，
支持背压（ready/valid握手协议）。

\subsection{验证结果}

\subsubsection{功能仿真验证}

使用ModelSim对全部RTL模块进行了功能仿真验证：

\begin{table}[htbp]
\centering
\caption{功能仿真测试结果}
\begin{tabular}{llcc}
\toprule
\textbf{测试项} & \textbf{测试内容} & \textbf{测试用例数} & \textbf{结果} \\
\midrule
稀疏感知分析器 & 全零/全非零/50\%/聚集度 & 4 & 全部通过 \\
粒度决策引擎 & 5种稀疏区间 + 强制模式 & 6 & 全部通过 \\
L2编码-译码往返 & 全零/全非零/对角/交替/随机 & 25 & 全部通过 \\
顶层集成 & 强制L2往返 + 自适应决策 & 2 & 全部通过 \\
\bottomrule
\end{tabular}
\end{table}

\subsubsection{关键测试用例说明}

\paragraph{稀疏率统计精度}全零张量稀疏率测量值256（Q8.8格式，理论值256），
误差为0；50\%稀疏张量测量值128，误差为0。

\paragraph{粒度决策正确性}90\%稀疏+低聚集→L1（reason=1），
90\%稀疏+高聚集→L2（reason=2），60\%稀疏→L2（reason=3），
35\%稀疏→L3（reason=4），20\%稀疏→dense（reason=5），
全部符合设计预期。

\paragraph{编码-译码往返正确性}25组测试向量（含5组定向+20组随机），
编码后译码恢复数据与原始数据完全一致，往返错误率为0。

\subsubsection{Python参考模型验证}

使用Python参考模型在真实模型维度上评估了稀疏编码的压缩效果：

\begin{table}[htbp]
\centering
\caption{不同稀疏率下的自适应编码压缩率}
\begin{tabular}{cccc}
\toprule
\textbf{稀疏率} & \textbf{自适应粒度} & \textbf{压缩率} & \textbf{节省} \\
\midrule
90\% & L1元素级 & 37.5\% & 62.5\% \\
80\% & L2块级 & 56.3\% & 43.7\% \\
70\% & L2块级 & 62.5\% & 37.5\% \\
60\% & L2块级 & 68.8\% & 31.2\% \\
50\% & L2块级 & 75.0\% & 25.0\% \\
40\% & L3通道级 & 85.0\% & 15.0\% \\
30\% & L3通道级 & 92.0\% & 8.0\% \\
20\% & dense回退 & 100\% & 0\% \\
\bottomrule
\end{tabular}
\end{table}

% ============================================================================
% 第三部分：关键保护点介绍
% ============================================================================
\section{关键保护点介绍}

\subsection{保护点一：运行时自适应多粒度决策方法}

\subsubsection{保护内容}

一种在AI芯片中根据运行时数据稀疏特性自适应选择编码粒度的方法，
其特征在于包括以下步骤：

\begin{enumerate}
  \item S1：对输入张量进行运行时稀疏特性分析，
        获取稀疏率$\rho_{\text{sparse}}$和空间聚集度$\rho_{\text{cluster}}$；

  \item S2：根据所述稀疏率和空间聚集度，从预设的多级稀疏粒度中自适应选择最优粒度，
        其中多级粒度至少包括元素级（L1）、块级（L2）和通道级（L3）三个级别，
        决策规则为：
        \begin{itemize}
          \item 若$\rho_{\text{sparse}} \geq T_{L1}$且$\rho_{\text{cluster}} < T_c$，
                选择元素级粒度；
          \item 若$\rho_{\text{sparse}} \geq T_{L1}$且$\rho_{\text{cluster}} \geq T_c$，
                或$T_{L2} \leq \rho_{\text{sparse}} < T_{L1}$，
                选择块级粒度；
          \item 若$T_{L3} \leq \rho_{\text{sparse}} < T_{L2}$，
                选择通道级粒度；
          \item 若$\rho_{\text{sparse}} < T_{L3}$，
                回退为稠密模式。
        \end{itemize}

  \item S3：按照所选粒度对输入张量进行稀疏编码，
        生成包含粒度标识的统一格式编码数据；

  \item S4：编码数据被送入计算单元时，
        根据粒度标识分发至对应的解码通路恢复非零数据参与计算。
\end{enumerate}

\subsubsection{保护理由}

该方法是本发明的核心创新，与现有技术的本质区别在于：

\begin{itemize}
  \item 现有技术（NVIDIA 2:4、SCNN等）均为编译时/训练时静态确定稀疏粒度；
  \item 本发明首次实现运行时基于数据特性的自适应粒度选择，
        决策延迟仅为若干时钟周期；
  \item 该方法使得同一IP核能够自动适配不同模型层的不同稀疏模式，
        无需人工调优。
\end{itemize}

\subsection{保护点二：统一混合粒度编码格式}

\subsubsection{保护内容}

一种支持多粒度混合的稀疏编码数据格式，其特征在于：

\begin{enumerate}
  \item 编码数据包包含2bit粒度标识字段，
        用于指示该包的编码粒度为元素级、块级、通道级或稠密中的至少一种；

  \item 粒度标识后的载荷格式根据粒度类型变化：
        \begin{itemize}
          \item 元素级（00）：载荷为坐标列表$\{(r_i, c_i)\}$与对应非零值$\{v_i\}$；
          \item 块级（01）：载荷为块位图$\text{bitmap}[B_H \times B_W]$与紧凑排列的非零值；
          \item 通道级（10）：载荷为通道掩码$\text{ch\_mask}[C]$与非零通道数据；
          \item 稠密（11）：载荷为原始稠密数据，无额外索引开销。
        \end{itemize}

  \item 同一数据流中的不同编码数据包可使用不同的粒度标识，
        译码器根据每个包的粒度标识独立解码。
\end{enumerate}

\subsubsection{保护理由}

\begin{itemize}
  \item 现有稀疏格式（CSC、CSR、COO、位图）均为单一粒度专用格式；
  \item 本发明的统一格式允许不同区域/层使用不同粒度，
        且格式切换仅增加2bit标识开销；
  \item 译码器无需预知粒度类型，从数据包中读取标识即可动态分发，
        实现零开销粒度切换。
\end{itemize}

\subsection{保护点三：稀疏感知分析器的硬件实现}

\subsubsection{保护内容}

一种用于运行时稀疏特性分析的硬件电路，其特征在于：

\begin{enumerate}
  \item 包含零值检测单元，对每个时钟周期输入的数据向量进行逐元素零值判断，
        输出非零掩码；

  \item 包含计数器阵列，实时累计总元素数、非零元素数和聚集块数，
        其中聚集块数的计算方法为：当当前周期存在非零元素且前一周期不存在非零元素时，
        聚集块计数器加1；

  \item 包含逐次逼近除法器，在张量处理结束时，
        计算$\rho_{\text{sparse}} = (N_{\text{total}} - N_{\text{nz}}) / N_{\text{total}}$
        的Q8.8定点表示，采用16步迭代完成，避免硬件除法器的高面积开销；

  \item 张量处理结束时输出统计结果，同时自动清零计数器，
        为下一个张量的统计做准备。
\end{enumerate}

\subsubsection{保护理由}

\begin{itemize}
  \item 现有稀疏加速器均不包含运行时稀疏特性分析硬件；
  \item 本发明首次实现硬件级实时稀疏率和聚集度统计，
        面积开销极小（仅需计数器和简单除法器）；
  \item 逐次逼近除法器相比直接硬件除法节省约70\%面积，
        16步迭代在200MHz时钟下仅需80ns。
\end{itemize}

\subsection{保护点四：零开销稠密回退机制}

\subsubsection{保护内容}

一种在稀疏编码IP核中实现零开销稠密回退的方法，其特征在于：

\begin{enumerate}
  \item 当稀疏率低于预设阈值（默认30\%）时，
        决策引擎自动选择稠密模式（granularity=11）；

  \item 稠密模式下，编码器直接透传原始数据，
        不进行任何编码操作，索引字段长度为0；

  \item 译码器根据粒度标识为11时，直接输出输入数据，
        无需经过任何解码逻辑；

  \item 稠密回退通路的延迟等于1个时钟周期（仅寄存器直通），
        与编码通路的选择完全并行，不影响关键路径时序。
\end{enumerate}

\subsubsection{保护理由}

\begin{itemize}
  \item 现有方案（如NVIDIA 2:4）强制要求50\%稀疏率，
        无法在低稀疏度时自动回退，导致低稀疏度场景性能倒退；
  \item 本发明的零开销回退机制确保IP核在任何稀疏度下都不会出现性能下降，
        最差情况等价于无稀疏加速的稠密计算。
\end{itemize}

\subsection{保护点五：可配置块级编码的可变块大小}

\subsubsection{保护内容}

一种在块级稀疏编码中支持可变块大小的方法，其特征在于：

\begin{enumerate}
  \item 块级编码器的块大小$B_H \times B_W$通过配置寄存器可编程设定，
        支持$2 \times 2$、$4 \times 4$、$8 \times 8$至少三种块大小；

  \item 块大小选择影响编码效率：
        \begin{itemize}
          \item 块越小（$2 \times 2$），位图开销越低，
                适合非零元素高度分散的场景；
          \item 块越大（$8 \times 8$），单个块可容纳更多非零元素，
                适合非零元素高度聚集的场景。
        \end{itemize}

  \item 块大小配置通过AXI-Lite寄存器cfg\_block\_size设定，
        可在运行时动态切换，无需重新综合。
\end{enumerate}

\subsubsection{保护理由}

\begin{itemize}
  \item 现有块级稀疏方案（如2:4）块大小固定，
        无法适配不同聚集度的稀疏分布；
  \item 本发明的可变块大小机制使得块级编码在不同聚集度下
        均能获得最优的编码效率。
\end{itemize}

\subsection{保护点六：粒度路由器的动态分发机制}

\subsubsection{保护内容}

一种根据编码数据中的粒度标识动态分发的路由器电路，其特征在于：

\begin{enumerate}
  \item 路由器接收包含2bit粒度标识的编码数据，
        根据标识值将数据分发至对应的编码/译码通路；

  \item 路由器为纯组合逻辑（1-to-4 demux），
        无流水线级，延迟为0个时钟周期；

  \item 四个输出通路（L1/L2/L3/dense）的valid/ready信号
        由路由器根据粒度标识选择性地传递，
        未选通通路的valid信号保持为0。

  \item 路由器支持背压传播：未选通通路的ready信号不影响输入侧，
        仅选通通路的ready信号回传至输入侧。
\end{enumerate}

\subsubsection{保护理由}

\begin{itemize}
  \item 现有方案无粒度路由需求（单一粒度），不存在类似机制；
  \item 本发明的零延迟路由器确保粒度切换不影响流水线吞吐，
        背压传播机制保证数据一致性。
\end{itemize}

\subsection{权利要求书框架}

基于上述保护点，建议的权利要求层级结构如下：

\begin{table}[htbp]
\centering
\caption{权利要求层级}
\begin{tabular}{clp{8cm}}
\toprule
\textbf{权利要求} & \textbf{类型} & \textbf{保护内容} \\
\midrule
1  & 独立-方法 & 运行时自适应多粒度决策方法（S1-S4） \\
2  & 独立-装置 & 自适应多粒度稀疏编码IP核 \\
3  & 从属 & 统一混合粒度编码格式 \\
4  & 从属 & 稀疏感知分析器的硬件实现 \\
5  & 从属 & 零开销稠密回退机制 \\
6  & 从属 & 可配置块级编码的可变块大小 \\
7  & 从属 & 粒度路由器的动态分发机制 \\
8  & 从属 & 决策阈值可编程配置 \\
9  & 从属 & 编码器阵列的并行流水线结构 \\
10 & 从属 & 与MAC阵列的耦合接口 \\
\bottomrule
\end{tabular}
\end{table}

\subsection{与现有技术的对比总结}

\begin{table}[htbp]
\centering
\caption{本发明与现有技术对比}
\begin{tabular}{lcccc}
\toprule
\textbf{对比维度} & \textbf{NVIDIA 2:4} & \textbf{EIE} & \textbf{SCNN} & \textbf{本发明} \\
\midrule
稀疏粒度     & 固定块级   & 固定元素级 & 固定元素级 & L1/L2/L3自适应 \\
决策时机     & 编译时静态 & 编译时静态 & 编译时静态 & 运行时动态 \\
编码格式     & 2:4专用    & CSC格式   & 坐标列表   & 统一混合格式 \\
适用模型     & 特定稀疏   & 非结构化  & 非结构化   & CNN/Transformer通用 \\
低稀疏回退   & 无         & 无        & 无        & 零开销稠密回退 \\
块大小可配   & 否         & 否        & 否        & 是（2/4/8） \\
\bottomrule
\end{tabular}
\end{table}

\vspace{2em}
\begin{center}
\textbf{--- 专利报告书完 ---}
\end{center}

\end{document}
